% STATUS: DRAFT
\chapter{Subsystems and entanglement}\label{ch:subsystems}

\begin{physmotiv}
In ordinary quantum mechanics a subsystem is a tensor factor, entanglement is a property of a pure state on $\HH_A\otimes\HH_B$, and the entropy of $\rho_A=\Tr_B\ket\psi\!\bra\psi$ measures it. In QFT none of this survives literally: the algebras of adjacent regions are not tensor factors. The algebraic replacement of ``subsystem'' is \emph{subalgebra}, and the replacements of ``trace over $B$'' are \emph{restriction of states} and \emph{conditional expectations}. This chapter develops those and states what remains true: the vacuum is entangled, in a precise and maximal sense, even though no entropy is defined.
\end{physmotiv}

\section{Subsystems as subalgebras}

For $\HH=\HH_A\otimes\HH_B$ the algebra of subsystem $A$ is $\M_A=\Bnd{\HH_A}\otimes\id$, with commutant $\M_A'=\id\otimes\Bnd{\HH_B}$. These are type I factors, and the pair $(\M_A,\M_A')$ has two special properties: (i) $\M_A\vee\M_A'=\BH$ (together they generate everything) and (ii) $\M_A\cap\M_A'=\C\id$. The algebraic definition:

\begin{definition}
A \emph{subsystem} of a quantum system with algebra $\M$ is a von Neumann subalgebra $\N\subset\M$. The ``complementary'' system is $\N'\cap\M$ (the part of $\M$ that commutes with $\N$).
\end{definition}

A state $\omega$ on $\M$ restricts to a state $\omega|_\N$ on $\N$: this is the algebraic partial trace. Even if $\omega$ is pure on $\M$, $\omega|_\N$ may be mixed: that is entanglement. In QFT the relevant data are the nets of algebras $\A(O)$ (\cref{ch:haagkastler}) and the restrictions of the vacuum.

\begin{whatwrong}
For two regions $O_1,O_2$ with disjoint closures, the algebra generated by $\A(O_1)$ and $\A(O_2)$ is isomorphic to $\A(O_1)\bar\otimes\A(O_2)$ only under the \emph{split property} (below). For \emph{adjacent} regions (sharing a boundary), such as the two Rindler wedges, no such isomorphism exists: $\A(R)\vee\A(L)$ is generated by two commuting type III$_1$ factors, but in a Hilbert space with a vacuum $\Omega$ the vector $\Omega$ is not a product vector in any tensor decomposition $\HH=\HH_R\otimes\HH_L$ compatible with the algebras.
\end{whatwrong}

\section{The split property}

In QFT, algebras of strictly nested regions are better behaved than adjacent ones.

\begin{definition}
A net of algebras has the \emph{split property} if for every pair of regions $O_1\Subset O_2$ ($\overline{O_1}\subset$ interior of $O_2$) there is a type I factor $\mathcal N$ with $\A(O_1)\subset\mathcal N\subset\A(O_2)$.
\end{definition}

Then $\HH\cong\HH_{\mathcal N}\otimes\HH_{\mathcal N'}$ and one may define a density matrix $\rho_{\mathcal N}$ and its entropy. The split property holds for free massive fields and, more generally, under a nuclearity condition on the energy dependence of local states (Buchholz--D'Antoni--Fredenhagen). It is the algebraic form of ``the vacuum has finite-energy-density correlations'': entanglement between regions separated by a gap decays fast enough to be captured by a trace-class map. The entropy of $\rho_{\mathcal N}$ blows up as the gap between $O_1$ and $O_2$ shrinks (the area-law divergence $\sim\mathrm{Area}/\epsilon^{d-2}$): the gap is the algebraic cutoff $\epsilon$. This makes ``entanglement entropy of a region'' a statement about the \emph{choice} of intermediate factor $\mathcal N$, not about $\A(O_1)$ alone.

\section{Conditional expectations: the algebraic partial trace}

For $\N\subset\M$, a \emph{conditional expectation} is a linear map $E:\M\to\N$ with $E(\id)=\id$, $E(n)=n$ for $n\in\N$, positivity, and the bimodule property $E(n_1mn_2)=n_1E(m)n_2$. It is the natural ``averaging over the complement''. For example, for $\M=\Bnd{\HH_A}\otimes\Bnd{\HH_B}$ and $\N=\Bnd{\HH_A}\otimes\id$, $E(a\otimes b)=a\,\omega_B(b)$ for a fixed state $\omega_B$ on $B$: ``take the expectation value in a state of $B$''. Expectations $E$ are not unique: they depend on a choice of reference state.

\begin{theorem}[Takesaki]
Let $\omega$ be a faithful normal state of $\M$ and $\N\subset\M$ a von Neumann subalgebra. There exists a normal conditional expectation $E:\M\to\N$ with $\omega\circ E=\omega$ iff $\N$ is invariant under the modular flow of $\omega$: $\sigma^\omega_t(\N)=\N$ for all $t$. The expectation is then unique.
\end{theorem}
This theorem (proved with the machinery of \cref{ch:tomita}) already says something physical: a subsystem can be ``traced out'' while preserving the state precisely when it is stable under the state's hidden dynamics. For a Rindler wedge in the vacuum, the modular flow is a boost (\cref{ch:rs_bw}), and subalgebras invariant under it are those generated by boost-invariant regions: this is why dilation-invariant subregions are the good ones (and the ones for which entropies are computed in 2D CFT via the replica trick). The \emph{Petz recovery map}, which undoes a quantum channel on a state when relative entropy is not decreased (\cref{ch:modham}), is built from the same data.

\section{Entanglement without entropy}

Even when no entropy exists we can still say: \emph{the vacuum is maximally entangled between $\A(O)$ and $\A(O)'$.}

\begin{theorem}[Summers--Werner and generalizations; hypotheses vary by version, see the literature]
For the vacuum of a QFT and the algebras $\M=\A(R)$, $\M'=\A(L)$ of complementary wedges (or of regions $O_1,O_2$ that touch or are only slightly separated), there are observables in $\M$ and $\M'$ for which the CHSH inequality is violated maximally, attaining the Tsirelson bound $2\sqrt2$. More generally, for type III factors, \emph{every} normal state violates Bell inequalities maximally for suitable pairs of commuting observables.
\end{theorem}
This is the sharpest algebraic version of the statement that the vacuum is entangled across any boundary: the maximal violation is a property of type III algebras and not of any special state. (It is also consistent with Reeh--Schlieder, \cref{ch:rs_bw}, which says that local operators acting on the vacuum can approximate any state.)

What \emph{is} well defined and finite in QFT: the relative entropy $S(\omega\|\varphi)$ of two states (Araki, \cref{ch:modham}), the mutual information $I(A:B)=S(\omega_{AB}\|\omega_A\otimes\omega_B)$ for regions with a gap, the \emph{differences} of entropies between states (``entanglement first law'' $\delta S=\delta\langle K\rangle$), and, with gravity, the generalized entropy (\cref{ch:clpw}).

\begin{keyidea}
Subsystems are subalgebras; entanglement means the restricted state is mixed or correlated in a non-product way; conditional expectations preserving a state exist iff the subalgebra is modular-flow-invariant; and in QFT the vacuum is as entangled as quantum mechanics permits without having a well-defined entropy.
\end{keyidea}

\section*{Exercises}
\begin{exercise}
For $\M=M_2\otimes M_2$ and $\N=M_2\otimes\id$, show that $E(a\otimes b)=a\,\omega_B(b)$ is a conditional expectation for any state $\omega_B$. Show $\omega_{AB}\circ E=\omega_{AB}$ iff $\omega_{AB}$ is a product $\omega_A\otimes\omega_B$ (with this $\omega_B$).
\end{exercise}
\begin{exercise}
Use Takesaki's theorem in finite dimensions: $\omega=\Tr(\rho\,\cdot)$ on $M_n$ and $\N$ a subalgebra. Show $\sigma^\omega_t(a)=\rho^{it}a\rho^{-it}$ and that $\sigma_t(\N)=\N$ iff $\N$ is invariant under conjugation by $\rho^{it}$. Find the $\omega$-preserving expectation for $\N=M_k\otimes\id_m$ and $\rho=\rho_A\otimes\rho_B$ and compare with partial trace.
\end{exercise}
\begin{exercise}
Compute the CHSH value $\langle A_0B_0+A_0B_1+A_1B_0-A_1B_1\rangle$ for the Bell state $(\ket{\uparrow\uparrow}+\ket{\downarrow\downarrow})/\sqrt2$ with spin observables in the $xz$-plane at the best angles, and show it equals $2\sqrt2$. Explain how the type III$_1$ theorem upgrades this from a single carefully prepared state to \emph{every} state.
\end{exercise}
\begin{exercise}
Heuristically estimate the entanglement entropy between $A(O_1)$, a ball of radius $r$, and the complement of a slightly larger ball of radius $r+\delta$ in a $(3+1)$-dimensional CFT-like theory using $S\sim c\,r^2/\delta^2$. Interpret $\delta$ as the algebraic cutoff in the split property and explain why the limit $\delta\to0$ does not exist.
\end{exercise}
